\SourceLeaf{274}
\[
\begin{aligned}
\frac{\mathrm{d}U}{\mathrm{d}w}&=\frac{\cos\varphi}{\cos\theta}\\
\frac{\mathrm{dd}U}{\mathrm{d}w^2}&=-\frac{3ee\cos\varphi}{\cos\theta^3}\mathord{.}cs\\
\frac{\mathrm{d}^3U}{\mathrm{d}w^3}&=\frac{ee\cos\varphi}{\cos\theta^5}(3cc-3ss+ee(12ccss+3s^4))
\end{aligned}
\]
\[
\frac{\mathrm{d}^4U}{\mathrm{d}w^4}=\frac{ee\cos\varphi}{\cos\theta^7}\mathord{.}cs(16-ee(49cc-13ss)-e^4(56ccss+29s^4))
\]
\[
\begin{aligned}
\frac{\mathrm{d}^5U}{\mathrm{d}w^5}={}&\frac{ee\cos\varphi}{\cos\theta^9}(-16cc+12ss+ee(49c^4-378ccss+9s^4)\\
&+e^4(628c^4ss+174ccs^4-54s^6)+e^6(268c^4s^4+220ccs^6+33s^8))
\end{aligned}
\]
\indent In this expansion of \(q\) in a series in \(p\), it is tacitly assumed that both quantities are expressed in parts of the arc equal in length to the radius. If, however, \(q\) is to be in seconds and \(p\) in degrees, the first term of the series must be multiplied by \(3600\), the second by \(\tfrac{3600\pi}{180}=20\pi\), the third by \(3600(\tfrac{\pi}{180})^2=\tfrac{1}{9}\pi\pi\), and so on. With this assumption, applying the formulas to our example gives the following numerical values. I write them in a form that avoids lengthy decimal fractions.
\[
\begin{aligned}
q={}&359556''69447\mathord{.}\frac{p}{100}\\
&+3041{,}386524\mathord{.}(\frac{p}{100})^2\\
&-946{,}260563\mathord{.}(\frac{p}{100})^3\\
&-4135{,}396057\mathord{.}(\frac{p}{100})^4\\
&+227{,}04342\mathord{.}(\frac{p}{100})^5
\end{aligned}
\]
Since \(p\) is to amount to only a few units in applications, this series always converges very quickly. To confirm the accuracy of these numbers, I carried out the calculation for \(p=-6\) and for \(p=+6\), that is, for
\[
\begin{aligned}
P+p&=46^0\ 42'\ 2''53251\quad\text{and for}\\
P+p&=58\quad42\quad2{,}53251
\end{aligned}
\]
using both the series and the finite formula 3. The series gives
